\section{Exact arithmetic integrals}
\label{appb-integrals}

This appendix follows source Appendix~B, pp.~26--28.  Its calculations
are \textbf{[Derived]} identities for the functions in Section~5;
their use as determinant bounds has the same local dependencies.

\subsection{Evaluation of the inner integral}
\label{appb-inner}

For $0\leq u<1$, set $d_0(u)=\min(u,1-u)$ and let $e(u)$ equal
$1$ for $u\leq1/2$ and $-1$ otherwise.  For $f=\{x\}$,
\[
 \ell(x,z)-2x=e(f)\bigl(\boldsymbol{1}_{z<d_0(f)}-2d_0(f)\bigr)
\]
away from a single irrelevant endpoint.  Hence source (B.1) remains
exactly applicable:
\begin{align}
 \int_0^{1/2}B^2\,dz&=d_0(g)(1-2d_0(g)),\notag\\
 \int_0^{1/2}AB\,dz&=e(f)e(g)
       \bigl(\min(d_0(f),d_0(g))-2d_0(f)d_0(g)\bigr).
 \label{appb-fractional-parts}
\end{align}
Together with \eqref{s5-tail-decomposition}, this is an explicit
rational formula for $R_9$ without a remaining $z$-integral.

Let $3=t_0<\cdots<t_{143}=20$ be the endpoints and the interior
points where $cx$ is integral for some
\[
 c\in\{2,2\alpha,2\lambda,2H,4\alpha,2(1-\alpha),2(1+\alpha)\}.
\]
These are the same breakpoints as source (B.2).  The changes of branch
in the minima and positive parts are included.  The quadratic
coefficient still vanishes because
\[
 2\lambda-12\lambda\alpha
 +2(H^2-H-\lambda^2)-18\alpha^2+6\alpha=0.
\]
Thus $R_9=a_i x+b_i$ on each open interval.  Taking its values at
$x_i=(2t_i+t_{i+1})/3$ and $y_i=(t_i+2t_{i+1})/3$ determines
$a_i=(R_9(y_i)-R_9(x_i))/(y_i-x_i)$ and
$b_i=R_9(x_i)-a_i x_i$.  The exact analogue of (B.3) is
\begin{equation}
 \int_3^{20}\frac{R_9(x)}{x^3}\,dx
 =\sum_{i=0}^{142}\left[
 a_i\left(\frac1{t_i}-\frac1{t_{i+1}}\right)
 +\frac{b_i}{2}\left(\frac1{t_i^2}-\frac1{t_{i+1}^2}\right)
 \right].
 \label{appb-inner-sum}
\end{equation}
Exact rational evaluation of the original floor-function expression
for $\Gamma_9$ gives \eqref{s5-finite-inner}; a third interior point
on each interval checks the affine interpolation.  There is also a
short independent comparison: $a_{9,i}=a_{5,i}+4\lambda$ and
$b_{9,i}=b_{5,i}$, so the increase is
\[
 4\lambda\left(\frac13-\frac1{20}\right)=\frac{629}{600}.
\]
In particular, every entry of source Table~3 has the explicit analogue
\[
 \int_j^{j+1}R_9(x)x^{-3}\,dx
 =\int_j^{j+1}R_5(x)x^{-3}\,dx+\frac{37}{10j(j+1)},
 \qquad j=3,\ldots,19.
\]
This gives all seventeen unit-interval values from one exact rule.

\subsection{Evaluation of the outer integral}
\label{appb-outer}

Write $T_9(y)=C_9(y)+S(y)$.  Resolving the minima, positive parts,
and floor in Section~5 gives the following replacement for source
Table~4; on each interval $T_9(y)=b+cy$.
\begin{center}
\begin{tabular}{rrrr}
\hline
$l$&$r$&$b$&$c$\\
\hline
$1/3$&$43/120$&$35/4$&$-8$\\
$43/120$&$37/100$&$76/5$&$-26$\\
$37/100$&$37/80$&$267/20$&$-21$\\
$37/80$&$1/2$&$119/20$&$-5$\\
$1/2$&$43/80$&$39/5$&$-5$\\
$43/80$&$37/60$&$57/4$&$-17$\\
$37/60$&$37/40$&$62/5$&$-14$\\
$37/40$&$1$&$5$&$-6$\\
$1$&$37/20$&$37/20$&$-1$\\
\hline
\end{tabular}
\end{center}
The old rank-related breakpoints $13/30$ and $13/20$ disappear.
Summing $b(r-l)+c(r^2-l^2)/2$ gives $305479/144000$.
Separately, the raw determinant integrals satisfy
\[
 \int C_9=\frac{14441}{7200},\qquad
 \int C_5=\frac{23329}{19200},\qquad
 \int(C_9-C_5)=\frac{45541}{57600},
\]
where all three integrals run from $1/3$ to $1$ and
$C_5=R_0-d$ is the source's cost.  The common scalar integral
$1851/16000$ provides an additional check on the comparison.

\subsection{The final rational margins}
\label{appb-margins}

\paragraph{[Blocked conclusion] The source's margins change sign.}
With the unchanged leading real coefficient
$U=-2733991/2000000$, the exact sums at the two source cutoffs are
\begin{align*}
 A_{9,200}+U
 &=\frac{189010431570697800693281}{94195881588024216000000}
   =2.006567892\ldots,\\
 A_{9,100000}+U
 &=\frac{11624665884437287177113589319057}
          {5887242599251513500000000000000}
   =1.974551870\ldots.
\end{align*}
Consequently the proposed decay coefficients $-1600(A_{9,M}+U)$
are negative.  In the exact format of source Appendix~B.3, the
differences from its target margins are
\begin{align*}
 -1600(A_{9,200}+U)-\frac{139}{5}
 &=-\frac{190647085013289721446281}{58872425992515135000},\\
 -1600(A_{9,100000}+U)-\frac{7907}{100}
 &=-\frac{11915606054639047909891714319057}
           {3679526624532195937500000000}.
\end{align*}
The limiting sum is also positive:
$A_{9,*}+U=1.974487121\ldots$.
These are failures of the required sign for this certificate.
A positive upper bound is not a lower bound on the actual polynomial
values, so none of these calculations asserts growth, rationality,
or impossibility of another construction.
